\uuid{BfUb}
\titre{Limite en un point}
\auteur{bodin ; hanani ; mzari}
\organisation{exo7}
\licence{CC BY-NC-SA 4.0}
\datecreate{2020-01-20}
\niveau{L1}
\module{Analyse}
\chapitre{Continuité, limite et étude de fonctions réelles}
\sousChapitre{Limite de fonctions}
\difficulte{2}
\theme{Limites des fonctions réelles, limite en un point}
\exo7qcmid{lille-1.limites.015}

\contenu{
\begin{question}
\qclassification{123.03}
\qtags{motcle=limite en un point}
Soit $f(x)= \frac{x}{|x-1|}-\frac{3x-1}{|x^2-1|}$. Quelles sont les assertions vraies ?
\begin{answers}

    \good{$\lim_{x\to 1} f(x)=0$}

    \bad{$\lim_{x\to 1} f(x)=1$}

    \bad{$f$ n'admet pas de limite en $-1$.}

    \good{$\lim_{x\to -1} f(x)=+\infty$}

\end{answers}
\begin{explanations}
En utilisant la définition de la valeur absolue,
 $f(x)=\left\{\begin{array}{ccc}\frac{x^2+4x-1}{1-x^2},& \mbox{si} \, x \le -1 \\
 \frac{1-x}{1+x} ,& \mbox{si} \,  -1\le x \le 1\\
   \frac{x-1}{1+x} ,& \mbox{si} \,  x \ge 1 \end{array}\right.$
\end{explanations}
\end{question}
}
